% Reproduces the figures of the Wikipedia article on Temperley--Lieb algebras.
\documentclass[11pt]{article}
\usepackage[margin=2.5cm]{geometry}
\usepackage{amsmath}
\usepackage{temperleylieb}
\usepackage{parskip}

\tlset{box=fill, gap=0.15, strand style={line width=1.3pt}}

\begin{document}

The generators of $TL_5(\delta)$ are:
\[ \TLgenerators{5} \]
From left to right, the unit $1$ and the generators $e_1, e_2, e_3, e_4$.

Multiplication on basis elements can be performed by concatenation: placing
two rectangles side by side, and replacing any closed loops by a factor
$\delta$, for example
$e_1e_4e_3e_2 \times e_2e_4e_3 = \delta\, e_1e_4e_3e_2e_4e_3$:
\[
  \TL{5}{e_1 e_4 e_3 e_2} \times \TL{5}{e_2 e_4 e_3}
  = \TL[boundaries]{5}{e_1 e_4 e_3 e_2 e_2 e_4 e_3}
  = \TL[reduce]{5}{e_1 e_4 e_3 e_2 e_2 e_4 e_3}
\]

The Jones relations can be seen graphically:
\begin{gather*}
  \TL{3}{e_1 e_1} = \delta\, \TL{3}{e_1} \\
  \TL{3}{e_1 e_2 e_1} = \TL{3}{e_1} \\
  \TL{4}{e_1 e_3} = \TL{4}{e_3 e_1}
\end{gather*}

The five basis elements of $TL_3(\delta)$ are the following:
\[ \TL[separate, reduce]{3}{1 e_2 e_1 {e_1 e_2} {e_2 e_1}} \]
From left to right, the unit $1$, the generators $e_2$, $e_1$, and
$e_1e_2$, $e_2e_1$.

The same, drawn vertically and with names:
\[ \TL[vertical, separate, reduce, names, gap=0.4]{3}{1 e_2 e_1 {e_1 e_2} {e_2 e_1}} \]

\end{document}
